\section{Prompt Templates}
\label{app:prompts}

\Cref{fig:prompt-templates} shows both prompt templates verbatim; they are included directly
from the experiment repository (\texttt{experiment/prompts/}). Placeholders in braces are
filled per task and language: \texttt{\{language\}} is \emph{Python}, \emph{Java}, or
\emph{C++}; \texttt{\{language\_note\}} is given in \cref{tab:language-notes};
\texttt{\{problem\}} is the problem statement (EffiBench-X) or the task specification
(CLBG); \texttt{\{starter\}} is the EffiBench-X code skeleton for the respective language.
\Cref{fig:prompt-example} shows a fully rendered prompt.

\begin{table}[h]
  \centering
  \caption{Language-specific line \texttt{\{language\_note\}}.}
  \label{tab:language-notes}
  \footnotesize
  \begin{tabular}{@{}lp{0.72\columnwidth}@{}}
    \toprule
    Language & Text \\
    \midrule
    Python & Target Python 3.14 (CPython). \\
    Java   & Target Java 25 (OpenJDK). \newline
             \emph{Program template only:} The public class must be named \texttt{Main}. \\
    C++    & Target C++20. The code is compiled with \texttt{g++ -O2 -std=c++20}. \\
    \bottomrule
  \end{tabular}
\end{table}

\begin{figure*}[t]
  \textbf{\footnotesize (a) Function template (EffiBench-X tasks)}
  \VerbatimInput[fontsize=\scriptsize,frame=single,framesep=2mm]{../experiment/prompts/template_function.txt}
  \vspace{1ex}
  \textbf{\footnotesize (b) Program template (CLBG tasks)}
  \VerbatimInput[fontsize=\scriptsize,frame=single,framesep=2mm]{../experiment/prompts/template_program.txt}
  \caption{Prompt templates used for all LLMs.}
  \label{fig:prompt-templates}
\end{figure*}

\begin{figure*}[t]
  \VerbatimInput[fontsize=\scriptsize,frame=single,framesep=2mm]{../experiment/prompts/rendered/EB03/cpp.txt}
  \caption{Fully rendered prompt for task EB03 (Sum of Digit Differences of All Pairs) in C++.}
  \label{fig:prompt-example}
\end{figure*}
